

The application of deep learning to subsurface modelling and reservoir simulation is a growing area of research with a high potential impact. By developing surrogate models using deep learning, simulations can be completed in considerably less time than conventional simulators that use numerical methods, once they are trained. In carbon capture and storage, surrogate modelling resulted in flow predictions being nearly 700,000 times quicker than existing methods \citep{gegewen_nested}. Research in geothermal and oil and gas simulations have shown similar results, reducing computational time by a factor of 100 and 1000 respectively \citep{DENG_geothermal, Huang_gnn_oil_gas}. These examples show the huge potential of using surrogate models, developed using deep learning, to reduce the computational time of reservoir simulations which enables rapid screening of potential sites and optimise operations. 

The storage of hydrogen is thought to play an important role in the transition to net zero, allowing seasonal energy storage at a grid scale \citep{krevor2023subsurface}. Excess energy produced by renewable sources, such as solar and wind, can be used to produce hydrogen through the electrolysis of water \citep{kumar2022overview}. This reduces the need for renewable energy curtailment, which led to the total curtailment of 10TWh of electricity in Great Britain alone during 2025 \citep{montel2025curtailment}. Not only is energy lost, but developers lose revenue, leading to extra costs for governments through subsidies. These secondary impacts can make investing in renewables, which is crucial for reaching net zero, unattractive for both private investment and government support \citep{iea2025renewables}. Underground hydrogen storage in depleted gas fields is an area of growing interest. This is due to the higher storage capacity compared to salt caverns, strong-understanding of the reservoir geology and low initial capital expenditure \citep{FU_uhs_Surrogate}. 

In studies focused on developing hydrogen storage surrogate models, synthetic datasets are generated due to the lack of active sites in the world. These datasets are then used to develop a surrogate model. These methods have proven to be highly successful at predicting hydrogen migration and reservoir plume evolution. \citet{FU_uhs_Surrogate} achieved high accuracy, achieving $R^2$ values of 0.996 and 0.991 for hydrogen saturation and pressure respectively. \citet{MAO_cushion} investigated the effects of cushion gas and developed a surrogate model. Hydrogen withdrawal efficiency, hydrogen purity, gas-water ratio and well injectivity were predicted with high accuracy. The results highlighted that cushion gas has negligible impacts on the performance metrics and storage performance improved with the number of cycles. In addition, the surrogate model allowed for the key performance metrics to be predicted five million times quicker than traditional physics simulators. 

Other approaches have focussed specifically on operational parameters. \citet{H2storeoptim} focussed on predicting hydrogen mass stored, gaseous hydrogen mass, hydrogen production rate and hydrogen dissolution rate. The surrogate model developed, only used the input rate to predict the focussed metrics. The study showed that long short-term memory models are the most effective for handling non-linear dynamic changes, which is most relevant for predicting production rate. 

 Although surrogate models allow for rapid prediction of saturation, pressure buildup and operational parameters, in areas such as carbon capture and storage and oil and gas production this is largely unnecessary. Both these fields are non-cyclical and therefore the models are often single-use. They are not designed to be continuously used or scaled. Unless the models developed can generalise with high accuracy to multiple reservoir sizes, heterogeneities and operational parameters they have limited use \citep{motaei2026advancements}. 

In systems which are cyclical like geothermal energy and gas storage, the use of surrogate models presents an area of research with a high potential impact. Due to their cyclical nature, similar scenarios can repeat multiple times through a storage sites life cycle. The primary issue is that hydrogen storage surrogate models currently use constant cycle lengths. In the studies by \citet{FU_uhs_Surrogate} and \citet{MAO_cushion}, one cycle consisted of six months of injection and six months of withdrawal. \citet{H2storeoptim} define a cycle as a five-month phase of hydrogen injection, a one-month shut-in phase and six-month production phase. Constant cycle structure is a limitation of the current studies \citep{mao2025deep}. Underground hydrogen storage is thought to operate in the same way as underground natural gas storage \citep{tarkowski2019underground}. The need for it derives from renewable energy curtailment. Existing natural gas storage sites do not operate with a constant cycle length and renewable energy curtailment varies seasonally \citep{pellow2015hydrogen}. Surrogate models for hydrogen storage must assess varying cycle lengths and rates of withdrawal and injection. Current hydrogen storage surrogate models also focus on aiming to generalise across multiple reservoir sizes and heterogeneities. For hydrogen storage in depleted gas fields this is largely unnecessary. Many of the sites that have been investigated for underground hydrogen storage (e.g. Roden, Viking A and Rough gas fields) have well-documented subsurface and formation understanding \citep{UHS_fields_Saadat}. As a consequence, surrogate models do not need to assess a large range of uncertainty in the heterogeneity of the fields; instead they should focus on operational parameters. 

In this study we demonstrate the use of surrogate modelling for hydrogen storage. We focus on assessing varying cycle lengths and operational rates which have been noted as an area for future work \citep{mao2025deep}. The cycle lengths and rates are generated using European natural gas storage sites data as an analogue. We predict injection and withdrawal rates, hydrogen stored, hydrogen saturation and pressure. We show our three dimensional surrogate model, developed using a U-Net architecture, to be highly accurate for predicting some of the key variables in hydrogen storage.